\chapter{The Other Hypothesis: The Case for NP = RP}
\label{ch:otherhyp}

Every chapter of Part IV so far has argued, or assumed, that the program fails on hard boards: that the level sets of a reduction board shatter to a depth growing with $N$, that the queens have the overlap gap property of Chapter~\ref{ch:ogp} in disguise, and that $\NP \ne \RP$. A book that tests a thesis should state the other hypothesis as its defenders would, and this chapter does. It makes the case that the queens, and by the main equivalence every problem in $\NP$, admit a polynomial time randomized sampler, that nothing measured in this book contradicts it, and that the arguments against it are arguments from belief. The authors do not hold this view. The case is presented as those who hold it would present it, with the counterarguments at the end.

\section{The argument from the data}

Begin with what has been measured, since the defenders of $\NP = \RP$ would begin there. Every quantity this book has computed on the queens appears polynomial at the sizes computed. The autocorrelation time of the bottom rung chain grows like $N^{1.8}$ to $N = 32$ with no curvature. The relaxation time, the worst case over observables, grows like $N^{3}$ to $N = 12$. The giant component of the bottom level set holds all but a vanishing fraction of solutions. The solutions are connected by near solution paths of length at most eight, and the maximum did not grow between $N = 11$ and $N = 13$. The second level set is connected at every $N$ examined. On conditioned boards the completions reconnect at two attacks and the level set at three or four, and the depth grew by one between $N = 12$ and $N = 16$, which is consistent with a logarithm. The repair neighbourhood is quadratic. The finite size correction to Simkin's theorem is zero. There is no thermal transition. Nowhere in a hundred pages of computation is there an exponential, and the defenders would say that a problem whose every measurable quantity is polynomial is a polynomial problem until shown otherwise.

The reply that the computations are on easy boards is anticipated. The defenders would answer that the boards the book calls easy are the boards that exist: random boards, planted boards, and the unconditioned board, which between them are every board a user of the queens will ever meet, and that the boards the book calls hard are a construction, the Gent Jefferson Nightingale encoding of a formula, which no one has built and whose level sets no one has looked at. The burden of proof, they would say, lies with the side that predicts an exponential it has never seen.

\section{The argument from structure}

The second argument is that the queens are not a random constraint satisfaction problem, and that the intuitions imported from one do not apply. The overlap gap property\index{overlap gap property} and the clustering transition of Chapter~\ref{ch:ogp} are theorems about random instances: random $k$ satisfiability, random regular graphs, the spin glass with Gaussian couplings. Their proofs use the independence of the constraints and the concentration that independence gives. The queens have no randomness in their constraints. The diagonals are the same on every board, the constraint graph is the same on every board, and the only randomness is in the fixed queens, which are few. A problem with a fixed, highly symmetric, long range constraint structure is closer to the problems where Monte Carlo succeeds than to the ones where it fails: the permanent, where every row constrains every column and Jerrum, Sinclair and Vigoda's chain mixes; the Ising model above its critical temperature, where the Swendsen Wang chain mixes; the volume of a convex body, where the constraints are global and the ball walk mixes. In each of these the constraints are dense and the level sets are connected, and in each the first intuition, that a dense constraint set must have a rugged landscape, was wrong.

The structural fact the defenders would point to is the one the book itself found: the solutions are not isolated. Table~\ref{tab:n1} of Chapter~\ref{ch:open} shows that most solutions have no solution neighbour under a single swap, and it was at first taken as the sign of a rugged landscape. The giant component result showed the opposite. The solutions are isolated from each other and connected through the sea of near solutions, which is exactly the geometry of the permanent, where perfect matchings are connected through near perfect matchings and the chain lives on the union. Jerrum, Sinclair and Vigoda's proof is a canonical path argument through that union, and the defenders would say that the queens want the same proof: a routing of each pair of solutions through near solutions, with the congestion bounded by an encoding argument of the kind that worked for matchings. The window divergences of Chapter~\ref{ch:open}, which say that the information in a solution is spread over half the board, are to them not an obstacle but the expected signature of a long range constraint, the same as for matchings, where the information in a perfect matching is spread over the whole graph.

\section{The argument from complexity theory}

The third argument is that nothing in complexity theory makes $\NP \ne \RP$ more than a belief. There is no proof that $\NP \ne \RP$, no proof that $\NP \ne \Pc$, and no proof that $\NP \not\subseteq \Ppoly$. The barriers of Chapter~\ref{ch:barriers} cut both ways: relativization, natural proofs and algebrization rule out the known techniques for separating the classes, and they say nothing in favour of the separation. The hardness versus randomness theory of Chapter~\ref{ch:pseudo} is a theory of $\Pc = \BPP$, which the defenders accept, and it says nothing about $\NP = \RP$; if anything, a world in which $\Pc = \BPP$ and $\NP = \RP$ is a world in which $\Pc = \NP$, which is Impagliazzo's Algorithmica, and Impagliazzo listed it as one of the five worlds because it has not been excluded.

The consequences of $\NP = \RP$ that Chapter~\ref{ch:consequences} calls implausible, the collapse of the polynomial hierarchy and the end of cryptography based on one way functions, are consequences, not contradictions. The polynomial hierarchy collapses under $\NP \subseteq \Ppoly$ as well, and no one has excluded that. One way functions are not known to exist, and Liu and Pass's theorem of Chapter~\ref{ch:kolmogorov} says they exist if and only if time bounded Kolmogorov complexity is hard on average, which is open. The defenders would add that the history of the subject is a history of surprises in the direction of tractability: linear programming, primality, the permanent for nonnegative matrices, the volume of convex bodies, graph isomorphism in quasipolynomial time, each believed hard by many until it was not. A separation that rests on the absence of algorithms, they would say, is only as strong as the search for algorithms has been thorough, and for sampling problems the search is forty years old.

\section{The argument from the equivalence}

The fourth argument turns the book's own theorems around. The main equivalence of Chapter~\ref{ch:theorems} says that $\NP = \RP$ if and only if the queens level sets on every board admit a polynomial sampler. The book reads this as a statement of what $\NP = \RP$ would cost. The defenders read it as a statement of what $\NP = \RP$ would take: one chain, or one generator, that mixes on one problem. The swap chain is one candidate, and its failure on a reduction board, if it fails, excludes one chain. The block chain of Chapter~\ref{ch:swendsen}, the cluster chain, the vertical walk with rung skipping, which does not need the bottom level set at all, the learned generator of Chapter~\ref{ch:gbc} with full history conditionals, and the countless chains no one has written down are the others. The rung skipping theorem of Chapter~\ref{ch:open} is, to them, the decisive move: the shattering of the bottom level set, which the book once took as the sign of hardness, costs a factor of thirty to forty at $N = 12$, and the question becomes whether the depth stays bounded, which the measurements to $N = 16$ do not exclude.

They would also note what the equivalence does not require. It does not require the sampler to be uniform; a polynomial bias suffices by Jerrum, Valiant and Vazirani. It does not require the sampler to work on every board in the same way; it may be adaptive, it may use the board's closure of Chapter~\ref{ch:sufficient} when the closure is large, and it may use the chain when it is small. And it does not require a proof of mixing; it requires mixing. The defenders would say that the program's honest position, that the chain works on every board anyone has tried and that its failure on untried boards is expected, is the position of a person who has found a polynomial algorithm and is reluctant to say so.

\section{The experiment the defenders would run}

The defenders would not argue; they would compute. Their experiment is the one the book has recommended since Chapter~\ref{ch:experiments2}: build the reduction board of a small hard formula, enumerate its level sets, and measure the shattering depth. If the depth is three or four, as on random boards, the hard boards are no harder than the random ones for the walk, and the case for $\NP = \RP$ has its first instance. If the depth is large, the case is in trouble on that board, and they would ask for the next chain. Their second experiment is the depth at $N = 20$ and $24$ on random boards, which separates $\log N$ from $N / 4$; their third is the Lanczos computation of the relaxation time to $N = 14$, which tells whether the exponent $3$ is stable; and their fourth is a learned generator on planted boards with the divergence reported, since a generator that reaches $\delta = O(\log N)$ on random boards at $N = 32$ would be the strongest evidence for the hypothesis that a computer can give.

\section{The reply}

The authors' reply is short, since the book is the long form of it. The data are polynomial because the boards are easy, and the boards are easy because they are random or planted; the sufficient statistic chapter argued that the natural training distribution avoids the hard boards, and the shattering depth, the one quantity measured on the hard side of the threshold, moves by about one between $N = 12$ and $16$, which the data do not yet show to be growth. The structural analogy with the permanent fails at the one place it must hold: perfect matchings of a bipartite graph are the common bases of two partition matroids, the near perfect matchings are the common independent sets one short of a basis, and the canonical path proof unwinds the symmetric difference of two matchings into alternating paths and cycles, a decomposition whose intermediate states stay within one defect, which the queens lack; two solutions cannot in general be interpolated by exchanges that stay within one attack, which is what the two step solution graph being nearly empty means. The complexity theoretic argument is correct as far as it goes, and it goes exactly as far as saying that $\NP \ne \RP$ is unproved, which the book has said on every page. And the equivalence argument is the one the authors find most serious, because it is true that the swap chain is one chain and that rung skipping removes the first obstacle; what remains is the depth, and a depth of $cN$ on reduction boards, which is what the encoding of a formula in the diagonals should produce, is an exponential that no change of chain removes, since every level set below it is disconnected for every local chain. The hypothesis that the depth is bounded on reduction boards is the hypothesis that a hard formula's satisfying assignments can be connected by local moves through placements with $O(1)$ violated clauses, and that is the statement the clustering theorems of Chapter~\ref{ch:ogp} deny for random formulas and the hard formula families deny by construction.

The two sides agree on the experiment. That is the state of the question.
